% Fonts & encoding
% Note: fontenc conflicts with fontspec in LuaLaTeX - removed for compatibility
% \usepackage[T1]{fontenc}  % Conflicts with fontspec
\usepackage{fontspec}
\usepackage{microtype}
% \usepackage{lmodern}  % Not needed with fontspec
\usepackage{enumitem}
% citations in chapters
\usepackage{epigraph}

% Define JuliaMono font for code listings (supports Unicode/Greek)
\newfontfamily\juliamono{JuliaMono}[
  Path=../assets/fonts/,
  Extension=.ttf,
  UprightFont=*-Regular,
  BoldFont=*-Bold,
  ItalicFont=*-LightItalic,
  BoldItalicFont=*-SemiBoldItalic
]



% Math & symbols
\usepackage{amsmath, amssymb, amsthm, mathtools, bm}

% Graphics & layout
\usepackage{graphicx}
\usepackage[margin=1in]{geometry}
\usepackage{booktabs}
\usepackage{tabularx}
\usepackage{threeparttable}
\usepackage{array}
\usepackage{makecell}
\newcolumntype{Y}{>{\raggedright\arraybackslash}X}
\usepackage{float}
\usepackage{caption}
\usepackage{subcaption}
\usepackage{chngcntr} % allow numbering floats by chapter


\usepackage[ruled,vlined,linesnumbered]{algorithm2e}
\usepackage{algpseudocode}
% \usepackage{algorithm}
% \usepackage[noend]{algpseudocode}
\SetKwInput{KwInput}{Input}
\SetKwInput{KwOutput}{Output}
\SetKw{KwReturn}{return}
\SetKw{KwInitialize}{initialize}
\SetKw{KwFunction}{function}
\SetKwComment{Comment}{}{}
% \usepackage[ruled,vlined,linesnumbered]{algorithm2e}
% \makeatletter
% \newcommand{\dontusealgpseudocode}{}
% \makeatotherx

% % Algorithm2e style settings
% \SetKwInput{KwInput}{Input}
% \SetKwInput{KwOutput}{Output}
% \SetKw{KwReturn}{return}
% \SetKw{KwInitialize}{initialize}
% \SetKw{KwFunction}{function}

\usepackage{etoolbox}
\BeforeBeginEnvironment{algorithm}{\par\vspace{0.5em}}
\AfterEndEnvironment{algorithm}{\par\vspace{0.5em}}


% Code listings (Julia-friendly) - using minted for better Unicode support
\usepackage{xcolor}
\definecolor{codebg}{RGB}{248,248,248}
\definecolor{dkgreen}{RGB}{0,100,0}
\definecolor{dkblue}{RGB}{0,0,139}
\definecolor{gray}{RGB}{102,102,102}
\definecolor{violet}{RGB}{148,0,211}
\definecolor{my_blue}{RGB}{0,114,178}
\definecolor{my_red}{RGB}{213,94,0}
\definecolor{my_green}{RGB}{0,158,115}


% \usepackage[cachedir=_minted-main]{minted}
% \usepackage[cache=true,cachedir=_minted-main]{minted}
\usepackage[cache=false]{minted}
\setminted{
  autogobble,
  breaklines=true,
  breakanywhere=true,
  breakautoindent=false,
  fontsize=\small,
  bgcolor=codebg,
  linenos=true,
  numbersep=8pt,
  frame=single,
  framesep=6pt,
  breaklines=true,
  tabsize=2
}



% % Enable page breaks within minted environments
\usepackage{etoolbox}
\AtBeginEnvironment{minted}{\par\vspace{0.5em}}
\AtEndEnvironment{minted}{\par\vspace{0.5em}}

% \renewenvironment{listing}{\begin{listing}[H]\begin{minipage}{\linewidth}}{\end{minipage}\end{listing}}

\usepackage{fvextra}
\DefineVerbatimEnvironment{JuliaREPL}{Verbatim}
{breaklines,fontsize=\small,formatcom=\color{black!90}}

\newcommand{\juliaout}[1]{\textcolor{gray!60}{#1}}

% Number listing floats per chapter: e.g., 2.1, 2.2, ...
\counterwithin{listing}{chapter}

% Use JuliaMono for minted code
%\setmonofont{JuliaMono}[
%  Path=../assets/fonts/,
%  Extension=.ttf,
%  UprightFont=*-Regular,
%  BoldFont=*-Bold,
%  ItalicFont=*-LightItalic,
%  BoldItalicFont=*-SemiBoldItalic
%]
% ADD this instead (uses fonts available on Overleaf):
\setmonofont{DejaVu Sans Mono}  % Good Unicode support, available on Overleaf
% \renewcommand{\ttdefault}{JuliaMono}

% Breakable, captioned code listings without floats
% (Removed temporary jucode helpers; using explicit captionof + minted instead)

% Ensure latexmk notices changes to external code files included via \inputminted
\newread\mintedtrackin
\newcommand{\mintedtrackline}{}
\newcommand{\mintedtrack}[1]{% reads one line to register file as INPUT in .fls
  \begingroup
  \openin\mintedtrackin=#1\relax
  \read\mintedtrackin to\mintedtrackline
  \closein\mintedtrackin
  \endgroup
  \par\ignorespaces
}

% Hyperlinks & clever references
% Define colors for links
\definecolor{webbrown}{rgb}{.6,0,0}  % Reddish-brown for citations
\usepackage{hyperref}
\usepackage{bookmark} % reduces rerun warnings for outlines
\hypersetup{
  breaklinks=true,
  colorlinks=true,
  linkcolor=black,
  citecolor=webbrown,
  urlcolor=webbrown,
  filecolor=black
}
\usepackage{cleveref}

% Avoid large vertical stretch/shrink warnings
\raggedbottom

% Theorems
\newtheorem{theorem}{Theorem}[chapter]
\newtheorem{proposition}[theorem]{Proposition}
\newtheorem{lemma}[theorem]{Lemma}
\theoremstyle{definition}
\newtheorem{definition}[theorem]{Definition}
\newtheorem{assumption}[theorem]{Assumption}
\theoremstyle{remark}
\newtheorem{remark}[theorem]{Remark}

% natbib (AER .bst expects author-year)
\usepackage[round,authoryear]{natbib}

% Callouts (tcolorbox)
\usepackage[most]{tcolorbox}
\tcbuselibrary{skins,breakable}

% (Reverted) Remove custom breakable minted tcolorbox code environments

% Optional icons for callout titles (fallbacks if fontawesome5 missing)
\IfFileExists{fontawesome5.sty}{
  \usepackage{fontawesome5}
  \newcommand{\iconnote}{\ifdefined\faIcon\faIcon{info-circle}\else\textbf{i}\fi}
  \newcommand{\icontip}{\ifdefined\faIcon\faIcon{lightbulb}\else\textbf{!}\fi}
  \newcommand{\iconimportant}{\ifdefined\faIcon\faIcon{exclamation-triangle}\else\textbf{!}\fi}
  \newcommand{\iconcaution}{\ifdefined\faIcon\faIcon{exclamation}\else\textbf{!}\fi}
}{
  \newcommand{\iconnote}{\textbf{i}}
  \newcommand{\icontip}{\textbf{!}}
  \newcommand{\iconimportant}{\textbf{!}}
  \newcommand{\iconcaution}{\textbf{!}}
}

% Callout color palette
\definecolor{calloutnote}{RGB}{30,100,180}
\definecolor{callouttip}{RGB}{0,130,90}
\definecolor{calloutimportant}{RGB}{180,60,0}
\definecolor{calloutcaution}{RGB}{160,0,0}


% Custom commands
\newcommand{\E}{\mathbb{E}}
\newcommand{\Var}{\mathrm{Var}}
\newcommand{\Cov}{\mathrm{Cov}}
\newcommand{\R}{\mathbb{R}}
\newcommand{\dd}{\mathrm{d}}

% Ensure figures go where asked (optional)
\usepackage{placeins}

% NOTE
\newenvironment{calloutnote}[1][Note]{%
  \begin{tcolorbox}[
    enhanced, breakable,
    boxrule=0.6pt, arc=2pt,
    left=8pt, right=8pt, top=6pt, bottom=6pt,
    coltitle=black, fonttitle=\bfseries,
    titlerule=0pt, toptitle=2pt, bottomtitle=2pt,
    borderline west={2.5pt}{0pt}{calloutnote},
    colback=calloutnote!6, colframe=calloutnote!55!black,
    colbacktitle=calloutnote!10,
    title={\small \iconnote\ \ #1}
  ]%
}{\end{tcolorbox}}

% TIP
\newenvironment{callouttip}[1][Tip]{%
  \begin{tcolorbox}[
    enhanced, breakable,
    boxrule=0.6pt, arc=2pt,
    left=8pt, right=8pt, top=6pt, bottom=6pt,
    coltitle=black, fonttitle=\bfseries,
    titlerule=0pt, toptitle=2pt, bottomtitle=2pt,
    borderline west={2.5pt}{0pt}{callouttip},
    colback=callouttip!6, colframe=callouttip!55!black,
    colbacktitle=callouttip!10,
    title={\small \icontip\ \ #1}
  ]%
}{\end{tcolorbox}}

% IMPORTANT
\newenvironment{calloutimportant}[1][Important]{%
  \begin{tcolorbox}[
    enhanced, breakable,
    boxrule=0.6pt, arc=2pt,
    left=8pt, right=8pt, top=6pt, bottom=6pt,
    coltitle=black, fonttitle=\bfseries,
    titlerule=0pt, toptitle=2pt, bottomtitle=2pt,
    borderline west={2.5pt}{0pt}{calloutimportant},
    colback=calloutimportant!6, colframe=calloutimportant!55!black,
    colbacktitle=calloutimportant!10,
    title={\small \iconimportant\ \ #1}
  ]%
}{\end{tcolorbox}}

% CAUTION
\newenvironment{calloutcaution}[1][Caution]{%
  \begin{tcolorbox}[
    enhanced, breakable,
    boxrule=0.6pt, arc=2pt,
    left=8pt, right=8pt, top=6pt, bottom=6pt,
    coltitle=black, fonttitle=\bfseries,
    titlerule=0pt, toptitle=2pt, bottomtitle=2pt,
    borderline west={2.5pt}{0pt}{calloutcaution},
    colback=calloutcaution!6, colframe=calloutcaution!55!black,
    colbacktitle=calloutcaution!10,
    title={\small \iconcaution\ \ #1}
  ]%
}{\end{tcolorbox}}    

\newcommand{\bE}{\mathbb{E}}
\newcommand{\bc}{\mathbf{c}}
\newcommand{\bej}{\mathbf{e}}
\newcommand{\ba}{\mathbf{a}}
\newcommand{\bs}{\mathbf{s}}
\newcommand{\bi}{\mathbf{i}}
\newcommand{\bj}{\mathbf{j}}
\newcommand{\bB}{\mathbf{B}}
\newcommand{\bP}{\mathbf{P}}
\newcommand{\bV}{\mathbf{V}}
\newcommand{\bI}{\mathbf{I}}
\newcommand{\bD}{\mathbf{D}}
\newcommand{\bff}{\mathbf{f}}
\newcommand{\bg}{\mathbf{g}}
\newcommand{\bU}{\mathbf{U}}
\newcommand{\bups}{\mathbf{\Upsilon}}
\newcommand{\bu}{\mathbf{u}}
\newcommand{\bv}{\mathbf{v}}
\newcommand{\br}{\mathbf{r}}
\newcommand{\bW}{\mathbf{W}}
\newcommand{\bZ}{\mathbf{Z}}
\newcommand{\bG}{\mathbf{G}}
\newcommand{\bH}{\mathbf{H}}
\newcommand{\bX}{\mathbf{X}}
\newcommand{\bY}{\mathbf{Y}}
\newcommand{\by}{\mathbf{y}}
\newcommand{\veta}{\boldsymbol{\eta}}
\newcommand{\bxi}{\boldsymbol{\xi}}
\newcommand{\bR}{\mathbf{R}}
\newcommand{\bx}{\mathbf{x}}
\newcommand{\bpi}{\bm{\pi}}
\DeclareMathOperator{\Tr}{Tr}
\newcommand{\hs}{\mathbf{H}}
\newcommand{\bt}{\mathbf{t}}
\newcommand{\bJ}{\mathbf{J}}
\newcommand{\hjb}{\textrm{HJB}}
\newcommand{\bepsilon}{\boldsymbol{\epsilon}}
\newcommand{\btheta}{\boldsymbol{\theta}}
\newcommand{\balpha}{\boldsymbol{\alpha}} 
\newcommand{\bsigma}{\boldsymbol{\sigma}} 
\newcommand{\bmu}{\boldsymbol{\mu}} 


\usepackage{tikz}
\usetikzlibrary{tikzmark}
\usetikzlibrary{calc,arrows.meta,positioning,tikzmark,backgrounds}
\usepackage{pgfplots}
\usepgfplotslibrary{patchplots,fillbetween}
\pgfplotsset{compat=1.18}
\tikzset{
  annarrow/.style={-{Latex[length=3mm]}, thick},
  anntext/.style={font=\small\bfseries}
}

% Node styles
\tikzset{
  neuron/.style     = {circle, draw=black!50, very thick, minimum size=8.5mm},
  inputnode/.style  = {neuron, fill=my_green!90},
  hiddennode/.style = {neuron, fill=my_blue!90},
  outputnode/.style = {neuron, fill=my_red!85},
  conn/.style       = {gray!70, -{Latex[length=2.2mm,width=1.4mm]}, line width=0.5pt},
  lbl/.style        = {font=\small},
  biglbl/.style     = {font=\normalsize\bfseries}
}

% ============================================================
% CES Notation and Conventions
% Include via \input{ces_notation} or \include{ces_notation}
% Assumes preamble with: amsmath, amsthm, tcolorbox, theorem environments
% ============================================================

% ------------------------------------------------------------
% Custom commands for CES notation (add to master preamble or keep here)
% ------------------------------------------------------------

% Aggregators (calligraphic)
\newcommand{\M}{\mathcal{M}}              % Generic mean
\newcommand{\Mbar}{\bar{\mathcal{M}}}     % Classical mean
\newcommand{\Ssum}{\mathcal{S}}           % Inner sum
\newcommand{\U}{\mathcal{U}}              % Utility over goods
\newcommand{\F}{\mathcal{F}}              % Production function
\newcommand{\Hfun}{\mathcal{H}}           % Output/allocation function
\newcommand{\V}{\mathcal{V}}              % Intertemporal utility
\newcommand{\CE}{\mathcal{R}}             % Certainty equivalent

% Weight vector
\newcommand{\bom}{\boldsymbol{\omega}}    % omega vector
\newcommand{\bell}{\boldsymbol{\ell}}     % ell vector

% Optimal values
\newcommand{\opt}{^{*}}                   % Superscript for optima

% Price index
\newcommand{\Pstar}{P^{*}}                % Ideal price index

% Expenditure shares
\newcommand{\wexp}{w}

\newcommand{\Rstar}{R^{*}}% Expenditure share

% KL divergence
\newcommand{\KL}[2]{D_{\mathrm{KL}}\left(#1 \,\|\, #2\right)}
